\chapter{水下光通信系统模型与理论基础}
\label{chap:theory_model}

可靠的水下无线光通信系统设计，需要以对水下光学传播机理和信道特性的准确认识为基础。本章围绕水下光学信道特性、时变双选择性信道模型、传统 OFDM 的局限性以及 AFDM 的基本原理展开分析，为后续传输方案设计奠定理论基础。

\section{水下光学信道特性分析}
\label{sec:underwater_channel}

光束在不同类型海水中传播时，会受到水分子、溶解物及悬浮颗粒的共同影响，主要表现为吸收、散射以及光学湍流等效应\cite{kaushalUnderwaterOpticalWireless2016a}。与陆地自由空间光通信相比，水下介质具有更强的吸收与散射特性，且其统计性质会随海域类型、颗粒浓度、温盐分布以及流体扰动状态的变化而明显改变。因此，UWOC 链路的传播损耗、可达距离以及时变特征都与具体水体环境密切相关。

从工程应用角度看，不同海域环境下的水下光链路往往表现出显著不同的传播条件。在清澈海水中，吸收与散射相对较弱，系统可获得更长的有效通信距离；而在近岸海域、港口水域或高浑浊水体中，悬浮颗粒浓度上升会显著增强散射效应，使链路衰减加剧并引入更明显的多径传播。由于本文关注的典型低动态场景往往对应浑浊水体中的短距链路，因此有必要从吸收、散射及其引发的时延扩展出发，对水下光信道的物理机理进行分析。

\subsection{吸收与散射效应}

吸收效应是指光子能量在与水体介质相互作用过程中转化为其他形式能量，从而导致光功率随传播距离增加而衰减。散射效应则是指光子在遇到悬浮颗粒后传播方向发生偏转。散射不仅会降低接收端收集到的有效光功率，还会因多重散射引起传播路径差异，进而造成时延扩展\cite{jaruwatanadilokUnderwaterWirelessOptical2008}。

在 UWOC 系统中，吸收与散射引起的综合衰减通常采用 Beer-Lambert 定律描述。设发射光功率为 $P_t$，传播距离为 $L$，则接收光功率 $P_r$ 可表示为
\begin{equation}
    P_{r} = P_{t} \exp(-c(\lambda) L)
\label{eq:beer_lambert}
\end{equation}
其中，$c(\lambda)$ 为总衰减系数，由吸收系数 $a(\lambda)$ 与散射系数 $b(\lambda)$ 组成，即 $c(\lambda)=a(\lambda)+b(\lambda)$\cite{zengSurveyUnderwaterOptical2017}。式(\ref{eq:beer_lambert}) 表明，在其他条件不变时，接收光功率随传播距离增加呈指数衰减。

需要指出的是，$a(\lambda)$ 与 $b(\lambda)$ 具有明显的波长相关性，因此波长选择会直接影响系统的传播损耗。现有研究表明，蓝绿光波段在海水中通常具有相对较低的总衰减系数，因此常被用作 UWOC 的主要工作窗口。这一波段能够在吸收与散射之间取得相对有利的平衡，使系统在有限发射功率下获得更大的传输距离和更高的接收信噪比。然而，不同水体环境对最优工作波长的影响并不完全一致：当水体浑浊程度提高时，散射增强会缩小系统的有效覆盖范围，并使链路性能对发射光束发散角、接收视场角及收发对准精度更加敏感。

除平均功率衰减外，散射过程还会对信号波形结构产生重要影响。对于非视距分量或多重散射分量，不同光子到达接收端的传播时间往往存在差异，因而会形成离散或连续分布的延迟分量。当信号带宽较高、符号周期较短时，这些延迟分量可能跨越多个采样间隔，从而使接收端观察到明显的时延扩展。时延扩展的直接后果是相邻符号之间发生叠加，并进一步导致频率选择性衰落。对于多载波系统而言，这意味着不同频率位置上的信道响应不再一致，部分子带可能出现深衰落，进而提高均衡与检测难度。

因此，在低动态浑浊水下场景中，影响系统性能的关键并不只是平均光功率损耗，还包括散射所引起的多径传播结构。换言之，即使收发平台相对运动较弱，信道仍可能由于介质本身的散射特性而表现出显著的频率选择性。这也是本文在低动态场景下重点关注多径鲁棒性设计的重要原因。

\subsection{水下光学湍流}

除吸收与散射外，温度梯度、盐度变化及洋流扰动还会导致水体折射率随机波动，形成水下光学湍流。湍流会引起波前畸变与光强闪烁，进而造成接收信噪比波动。在理论建模中，弱湍流条件下常采用对数正态分布描述衰落，中强湍流条件下则可采用 Gamma-Gamma 分布或其他复合统计模型进行表征\cite{zediniEvaluationUnderwaterOptical2025, huPerformanceCoherentOptical2025}。

与大气光通信中的湍流效应类似，水下光学湍流同样会导致接收功率的随机起伏，但其形成机理更容易受到局部温盐分层、流场扰动以及设备运行热源的影响。在短距离 UWOC 链路中，湍流造成的平均衰减未必占主导地位，但其对瞬时信噪比波动、误码性能离散性以及链路稳定性的影响仍不可忽略。当系统采用高阶调制或高灵敏度接收架构时，湍流引起的随机功率波动还可能放大同步误差与信道估计误差对检测性能的影响。

对于高动态场景而言，湍流效应还可能与平台运动造成的动态指向误差共同作用。平台姿态变化、机械振动以及波浪扰动会使收发端光束轴线发生偏移，导致接收面上的能量分布变化。当接收孔径较小或视场角较窄时，这类指向误差会进一步放大链路增益波动。尽管本文在后续模型中主要聚焦多径、多普勒和相位噪声影响，但从物理层角度看，湍流与指向误差仍是影响真实 UWOC 系统稳健性的重要因素，也是未来实验验证与模型拓展中需要进一步考虑的内容。

综上，水下光信道并非单一的功率衰减通道，而是同时包含吸收、散射、湍流及动态指向误差等多种扰动机制的复合信道。对于低动态场景，散射主导的时延扩展是影响系统设计的关键；对于高动态场景，平台运动与器件非理想性带来的时变效应则更加突出。对这些传播特性的分析，为后续建立双选择性信道模型和设计针对性传输方案提供了物理依据。

\section{时变双选择性信道数学模型}
\label{sec:doubly_selective_model}

在高动态 UWOC 场景中，水体散射引起的多径扩展与平台相对运动引起的多普勒扩展共同作用，使信道呈现明显的时变双选择性特征\cite{bemaniAFDMFullDiversity2021}。所谓双选择性，是指信道既在频率维度上表现出选择性衰落，又在时间维度上呈现快速变化特性。前者主要由时延扩展造成，后者主要由多普勒扩展引起。对于面向高速传输的 UWOC 系统，这两类效应往往难以彼此分离处理，因此需要统一的数学模型进行描述。

\subsection{多径与多普勒效应}

设最大时延扩展为 $\tau_{\max}$，最大多普勒频移为 $\nu_{\max}$。当时延扩展相对于符号周期不可忽略时，系统将表现出频率选择性衰落；当多普勒频移导致信道在一个传输帧内发生明显变化时，系统将表现出时间选择性衰落。这两类效应共同作用，会显著增加系统均衡与检测的难度。

从系统分析角度看，时延扩展通常可与相干带宽联系起来理解。若信号带宽超过信道相干带宽，不同频率分量将在传播过程中经历不同幅相响应，系统便会表现出明显的频率选择性。在多载波系统中，这意味着不同子载波所经历的信道增益不再一致，部分子载波甚至可能落入深衰落区。类似地，多普勒扩展通常可与相干时间相关联。当传输帧持续时间接近或超过信道相干时间时，信道在一帧内部便不能再视为准静态，从而使基于固定子载波正交关系或固定均衡系数的设计难以保持有效。

对于 UWOC 而言，低动态场景与高动态场景的主导失真机制并不相同。在低动态浑浊水体中，多径扩展往往更突出，系统的主要困难在于如何抑制时延扩展带来的频率选择性衰落与符号间干扰；而在高速移动场景中，平台相对运动会显著增强多普勒扩展，使时间选择性和载波间耦合问题变得更加突出。当两类效应同时存在时，系统接收端所面对的便不再是单纯的卷积信道，而是一个具有时变耦合结构的双选择性信道。

\subsection{基带等效信道表示}

对于包含 $P$ 条离散路径的水下线性时变信道，其时变基带等效冲激响应可写为
\begin{equation}
    h(t, \tau) = \sum_{i=1}^{P} h_{i} e^{j 2 \pi \nu_{i} t} \delta(\tau - \tau_{i})
\label{eq:ltv_channel}
\end{equation}
其中，$h_i$、$\tau_i$ 与 $\nu_i$ 分别表示第 $i$ 条路径的复增益、传播时延和多普勒频移。在离散时间域中，若采样间隔为 $T_s$、系统维度为 $N$，则该信道可等效表示为一个 $N \times N$ 的时变矩阵。由式(\ref{eq:ltv_channel}) 可见，不同路径同时引入时延与多普勒扩展，因而难以借助传统傅里叶域表示实现有效对角化。

进一步地，若令离散时延为 $l_i=\tau_i/T_s$，归一化多普勒为 $\alpha_i=\nu_i N T_s$，则每一路径均可视为“循环移位 + 相位旋转”的组合操作。于是，整个信道可以看作多个时移--频移算子的叠加。对于仅包含时延扩展而无明显多普勒扩展的情况，信道矩阵近似为循环卷积结构，较易通过 OFDM 等体制实现对角化处理；但当多普勒项不可忽略时，信道矩阵将由近似对角形式演化为具有带状耦合特征的时变矩阵，子载波之间会出现明显的能量泄露与相互干扰。

从连续模型过渡到离散模型的意义在于，后续系统设计和算法推导通常都以离散基带信号为研究对象。通过将信道表示为有限维矩阵形式，可以更方便地讨论不同波形在等效变换域中的结构特征，并据此构造检测与均衡方法。特别是对于本文关注的 AFDM 体制，信道在合适变换域中的耦合结构将直接决定后续参数设计和检测算法的复杂度。

\subsection{双选择性信道对系统设计的影响}

双选择性信道对通信系统设计的影响主要体现在三个方面。其一，在波形层面，系统需要同时应对时延扩展与多普勒扩展，而非仅针对单一失真机制进行优化。若波形仅在频域维持严格正交关系，则当多普勒扩展增强时，正交结构容易被破坏；若波形只能处理时域色散，则又难以兼顾快速时变引起的频域耦合。

其二，在接收算法层面，双选择性信道会使检测问题从逐子载波独立判决演化为带有耦合约束的联合估计问题。对于传统 OFDM，当信道近似静止时，各子载波可视为独立衰落；但在时变场景中，ICI 的出现使接收端需要考虑邻近甚至更大范围子载波之间的干扰耦合。若再叠加相位噪声、定时偏移或硬件非理想因素，检测复杂度将进一步上升。

其三，在系统实现层面，双选择性信道还会提高参数设计与同步处理的敏感性。例如，循环前缀长度、子载波间隔、帧长以及导频结构等参数都会对系统在时延扩展与多普勒扩展下的鲁棒性产生影响。对于 UWOC 这类短距高速链路而言，符号持续时间较短、系统带宽较大，因此即使绝对传播距离有限，仍可能对时延扩展和时变效应表现出较高敏感性。这也是本文后续分别针对低动态和高动态场景设计不同 AFDM 适配方案的原因。

\section{传统 OFDM 系统原理及其在时变信道中的局限性}
\label{sec:ofdm_limitations}

OFDM 通过离散傅里叶变换将高速串行数据映射到多个正交子载波上，从而将频率选择性信道分解为并行子信道\cite{armstrongOFDMOpticalCommunications2009}。在低动态场景下，若循环前缀长度足够，OFDM 可以有效抑制由时延扩展引起的符号间干扰。

从基本原理上看，OFDM 的核心思想是通过并行化传输将宽带频率选择性信道转化为多个窄带近似平坦衰落子信道。若循环前缀长度大于信道最大时延扩展，则线性卷积可近似转化为循环卷积，接收端在频域上能够实现逐子载波的一拍均衡。这一机制使 OFDM 在准静态多径环境中具有较高的实现效率和较成熟的工程基础，因此被广泛应用于射频通信和光通信系统中。

然而，OFDM 的有效性依赖于两个重要前提：一是循环前缀足以覆盖主要时延扩展，二是信道在一个 OFDM 符号周期内近似保持不变。当前提条件成立时，各子载波之间的正交关系可以较好维持，系统可将均衡与检测问题分解为多个独立的标量问题；而一旦信道在符号内部快速变化，这种独立性就会被破坏。

然而，在时变双选择性信道中，OFDM 对子载波正交性的依赖会成为性能瓶颈。当存在多普勒频移时，子载波之间会发生能量泄露并形成载波间干扰\cite{zhuDopplerResistantOrthogonalChirp2022}。此外，在深度频率选择性衰落条件下，OFDM 对局部频率位置的信道退化较为敏感，因此在高动态场景中难以保持稳定性能。

具体而言，多普勒频移会使接收信号相对于理想子载波基函数产生额外的相位旋转和频率偏移，从而破坏子载波间的严格正交关系。其结果是，原本应集中在某一子载波上的能量会扩散至相邻子载波，形成 ICI。随着多普勒扩展增加，ICI 的影响范围和强度都会上升，接收端不再能够通过简单的一拍均衡恢复发送符号，而需要引入更复杂的频域干扰抵消或时变信道估计方法。

在 UWOC 场景中，OFDM 还面临来自光学发射与接收机制的额外约束。对于 IM/DD 系统，发射信号需要满足实值与非负约束，因此通常需要通过厄米特对称、直流偏置或限幅等手段将复数 OFDM 基带信号转换为可驱动光源的单极性信号。虽然 DCO-OFDM、ACO-OFDM 等方案在满足非负约束方面各有优势，但这些处理方式也会带来频谱利用率下降、限幅噪声引入或平均功率开销增加等问题。当信道进一步呈现明显多径或时变特性时，传统光 OFDM 的优势将被削弱。

对于高动态场景，即使系统采用相干接收并恢复复数域信号，OFDM 仍难以从根本上摆脱对频域正交结构的依赖。尤其在多径、多普勒与相位噪声共同存在时，OFDM 接收端面对的是一个同时具有频域耦合与相位扰动的检测问题，系统性能和实现复杂度都会受到明显影响。因此，若要同时兼顾低动态多径场景和高动态时变场景的传输需求，仅对传统 OFDM 进行局部修正往往难以达到理想效果，仍有必要探索更适配时延--多普勒耦合信道的新型波形。

\section{AFDM 技术原理}
\label{sec:afdm_theory}

针对上述传统 OFDM 在双选择性信道中的局限，本节介绍 AFDM 的核心技术原理。AFDM 基于 DAFT 构造信号，通过引入啁啾参数改变信号在时频平面中的分布形式，使其在面对时延与多普勒共同作用时，能够在变换域中获得更具规律性的等效表示\cite{bemaniAFDMFullDiversity2021}。以下依次介绍 DAFT 的数学定义、啁啾参数设计原则及 AFDM 对本文两类场景的支撑作用。

\subsection{DAFT}

AFDM 的调制与解调分别通过离散逆仿射傅里叶变换（Inverse Discrete Affine Fourier Transform，IDAFT）与 DAFT 实现。其离散基带信号可写为
\begin{equation}
    s[n] = \frac{1}{\sqrt{N}} \sum_{k=0}^{N-1} X[k] e^{j 2 \pi (c_1 n^2 + \frac{nk}{N} + c_2 k^2)}
\label{eq:afdm_mod}
\end{equation}
其中，$X[k]$ 为第 $k$ 个 DAFT 域位置上的复数调制符号，$N$ 为总子载波数，$c_1$ 与 $c_2$ 分别为发射端与接收端啁啾参数。参数 $c_1$ 控制时域信号的二次相位调频速率，决定了啁啾基函数在时频平面中的瞬时频率变化轨迹；参数 $c_2$ 则作用于频域符号的预编码相位，影响不同子载波之间的相位耦合关系。当 $c_1 = c_2 = 0$ 时，式~\eqref{eq:afdm_mod} 退化为标准的离散傅里叶逆变换（IDFT），即 AFDM 退化为传统 OFDM。

上述调制过程可进一步写为矩阵形式。定义发射端啁啾相位对角矩阵 $\boldsymbol{\Lambda}_{c_1} = \mathrm{diag}\!\bigl(e^{j2\pi c_1 n^2}\bigr)_{n=0}^{N-1}$，接收端啁啾相位对角矩阵 $\boldsymbol{\Lambda}_{c_2} = \mathrm{diag}\!\bigl(e^{j2\pi c_2 k^2}\bigr)_{k=0}^{N-1}$，以及 $N$ 点归一化 IDFT 矩阵 $\mathbf{F}^{H}$，则 IDAFT 调制可以表示为
\begin{equation}
    \mathbf{s} = \boldsymbol{\Lambda}_{c_1} \mathbf{F}^{H} \boldsymbol{\Lambda}_{c_2} \mathbf{X}
\label{eq:afdm_matrix}
\end{equation}
其中，$\mathbf{X} = [X[0], X[1], \ldots, X[N-1]]^{T}$ 为 DAFT 域符号向量，$\mathbf{s} = [s[0], s[1], \ldots, s[N-1]]^{T}$ 为时域发射信号向量。式~\eqref{eq:afdm_matrix} 表明，AFDM 调制可以分解为三步级联操作：首先由 $\boldsymbol{\Lambda}_{c_2}$ 对频域符号施加二次相位预编码，然后经 IDFT 变换至时域，最后由 $\boldsymbol{\Lambda}_{c_1}$ 施加时域啁啾相位调制。相应地，接收端的 DAFT 解调为上述过程的逆操作。这种分解结构不仅有利于利用 FFT 实现高效调制与解调，也使得时延与多普勒对等效信道矩阵的影响可以在统一的矩阵框架下进行分析。

对于本文关注的两个典型场景，DAFT 的意义也有所不同。在低动态浑浊多径场景中，DAFT 的结构优势主要体现在对多径分量的分离能力以及对频率选择性衰落的适应性上；而在高动态场景中，DAFT 则有助于将多径与多普勒扩展映射为更具规律性的耦合形式，为时变信道下的检测与均衡提供便利。因此，DAFT 不仅是 AFDM 的数学工具，也是连接本文两类场景设计的统一理论基础。

\subsection{啁啾参数设计与全分集增益}

AFDM 的一个主要特点在于其可通过合理参数设计提升系统对双选择性信道的适应能力。根据相关研究，当参数满足特定条件时，AFDM 可使离散线性时变信道在 DAFT 域中呈现更具结构性的表示形式，从而有利于后续均衡与检测\cite{bemaniAFDMFullDiversity2021, taoAffineFrequencyDivision2025}。

从物理意义上看，参数 $c_1$ 与 $c_2$ 决定了发射与接收啁啾基函数在时频平面中的弯曲程度和相对匹配方式。不同参数取值会改变信号对时延与多普勒扰动的敏感性，也会影响等效信道矩阵在变换域中的稀疏程度和局部耦合结构。若参数设计合理，信道中不同路径对应的能量可在 DAFT 域中形成较为分离或可控的分布，从而提高系统对双选择性信道的适应能力；反之，若参数配置不当，不同路径或不同多普勒分量的响应可能发生严重重叠，降低后续检测性能。

所谓全分集增益，本质上是指系统能够充分利用多径信道中不同传播分量所蕴含的分集信息。对于传统 OFDM，由于信号主要在固定频率基上展开，当存在快速时变或深度选择性衰落时，部分子载波上的性能劣化会直接制约整体检测效果。AFDM 则通过啁啾映射改变了信息在时频域中的分布方式，使符号在传播过程中能够以更结构化的形式感知多个时延和多普勒分量，从而在合适条件下获得更好的分集利用能力。

需要强调的是，AFDM 的参数设计并非脱离具体场景而独立存在。对于本文的低动态场景，参数选择重点在于兼顾 IM/DD 架构下变换实现的稳定性与多径分离能力；对于高动态场景，参数设计则更关注时延--多普勒耦合信道下等效矩阵的局部结构，以便支撑后续检测算法设计。也就是说，AFDM 作为统一主线虽然贯穿全文，但其参数配置和具体实现形式仍需要根据场景特征进行适配。

\subsection{AFDM 对本文两类场景的支撑作用}

综合前述分析可以看出，AFDM 的核心价值在于为双选择性或近双选择性信道提供一种更具结构性的表示与处理方式。对于低动态浑浊多径场景，虽然多普勒扩展相对较弱，但由散射引起的多径传播仍会导致显著的频率选择性衰落。此时，若将 AFDM 与满足 IM/DD 非负约束的 ACO 机制结合，可在保持光学发射可实现性的同时增强系统对多径扩展的适应能力。对于高动态移动场景，AFDM 则能够与相干接收、索引调制及结构化检测方法进一步结合，以提高系统对多普勒扩展和相位噪声共同作用的鲁棒性。

因此，本章所讨论的水下光传播机理、双选择性信道模型、传统 OFDM 的局限以及 AFDM 的基本原理，共同构成了后续两章方案设计的理论基础。第三章将在 IM/DD 架构下研究面向低动态浑浊多径场景的 ACO-AFDM 传输方案；第四章则将在相干接收架构下研究面向高动态场景的 AFDM-IM 传输方案，并进一步讨论相位噪声与检测算法问题。
